\chapter{2015:Tomas Lindahl and Paul Modrich and Aziz Sancar}

\label{chap:chemistry2015}

The Nobel Prize in Chemistry for the year 2015 was awarded jointly to Tomas Lindahl and Paul Modrich and Aziz Sancar.

\textbf{Motivation:}

for mechanistic studies of DNA repair

\section{Tomas Lindahl}

\subsection*{Early Life and Education}

Tomas Lindahl was born in 1938 in Stockholm, Sweden.

\subsection*{Career and Major Research}

Lindahl received his doctorate from the Karolinska Institutet in Stockholm in 1967. He later moved to Britain, where he led a laboratory at the MRC laboratories at Clare Hall and became an emeritus scientist at the Francis Crick Institute in London. He demonstrated that DNA is chemically unstable and identified the repair enzymes that remove damaged bases.

\subsection*{Nobel Prize-winning Work}

The work for which Tomas Lindahl was awarded the Nobel Prize in Chemistry in 2015 was described by the Nobel Committee as: "for mechanistic studies of DNA repair".

Lindahl showed that bases are lost or damaged in DNA even without external agents, and he discovered the base excision repair pathway that removes and replaces such lesions.

\subsection*{Significance and Impact}

His discovery of base excision repair established that cells constantly protect their DNA against spontaneous damage, which would otherwise cause mutations and cancer.

\subsection*{Later Career and Honors}

Lindahl was director of the Cancer Research UK London Research Institute at Clare Hall and is now an emeritus group leader at the Francis Crick Institute.

\subsection*{Legacy}

Base excision repair is recognized as one of the essential safeguards of genetic information in all living cells.

\subsection*{Modern Developments}

Lindahl's discovery of base excision repair revealed the fundamental mechanism by which cells repair the constant chemical damage to DNA. His work established that DNA repair systems maintain the integrity of the genome and underpins modern cancer research, since defects in repair pathways are linked to hereditary cancers and drive many targeted cancer therapies.

\subsection*{Selected Publications}

"An N-Glycosidase from Escherichia coli That Releases Free Uracil from DNA Containing Deaminated Cytosine Residues" (Proceedings of the National Academy of Sciences, 1974)

\clearpage

\section{Paul Modrich}

\subsection*{Early Life and Education}

Paul Modrich was born in 1946 in Raton, New Mexico, USA.

\subsection*{Career and Major Research}

Modrich earned a bachelor's degree from the Massachusetts Institute of Technology in 1968 and a doctorate from Stanford University in 1973. He became a professor of biochemistry at Duke University and an investigator of the Howard Hughes Medical Institute, where he worked out the mechanism of mismatch repair.

\subsection*{Nobel Prize-winning Work}

The work for which Paul Modrich was awarded the Nobel Prize in Chemistry in 2015 was described by the Nobel Committee as: "for mechanistic studies of DNA repair".

Modrich identified the proteins that recognize and correct errors made when DNA is copied, and he reconstructed the mismatch repair process in detail.

\subsection*{Significance and Impact}

Mismatch repair catches the mistakes introduced during DNA replication, reducing the mutation rate enormously. Defects in this system predispose people to hereditary colorectal cancer.

\subsection*{Later Career and Honors}

Modrich continues his research on DNA mismatch repair at Duke University.

\subsection*{Legacy}

His work established mismatch repair as a central guardian of the genome and linked its failure to human disease.

\subsection*{Modern Developments}

Modrich's elucidation of mismatch repair, the DNA repair pathway that corrects replication errors, revealed how cells achieve the astonishing fidelity of DNA replication. His work explains the hereditary colon cancer syndrome Lynch syndrome and underpins modern cancer genetics and the development of immunotherapy for mismatch repair-deficient tumors.

\subsection*{Selected Publications}

"DNA Mismatch Correction" (Annual Review of Biochemistry, 1987)

\clearpage

\section{Aziz Sancar}

\subsection*{Early Life and Education}

Aziz Sancar was born in 1946 in Savur, Turkey.

\subsection*{Career and Major Research}

Sancar earned a medical degree from Istanbul University and a doctorate in molecular biology from the University of Texas at Dallas in 1977. He became a professor at the University of North Carolina at Chapel Hill, where he characterized the nucleotide excision repair system.

\subsection*{Nobel Prize-winning Work}

The work for which Aziz Sancar was awarded the Nobel Prize in Chemistry in 2015 was described by the Nobel Committee as: "for mechanistic studies of DNA repair".

Sancar identified the UvrABC enzyme system that cuts damaged DNA on both sides of a lesion and showed how the damaged piece is removed and replaced, and he studied the human equivalent of this process.

\subsection*{Significance and Impact}

Nucleotide excision repair removes damage caused by ultraviolet light and many chemicals. Sancar's work explained how cells remove such bulky lesions and why defects in the system cause skin cancer.

\subsection*{Later Career and Honors}

Sancar continues his research on DNA repair and its links to cancer at the University of North Carolina at Chapel Hill.

\subsection*{Legacy}

The UvrABC system he characterized is a paradigm for repair processes, and his studies connect DNA repair to cancer and chemotherapy.

\subsection*{Modern Developments}

Sancar's mapping of nucleotide excision repair, the pathway that repairs DNA damage caused by ultraviolet light, explained how cells protect the genome from sunlight and environmental mutagens. His work revealed the molecular clockwork of DNA repair and underpins the understanding of skin cancer risk and the design of chemoprevention strategies.

\subsection*{Selected Publications}

"DNA Excision Repair" (Annual Review of Biochemistry, 1996)

\clearpage

\section*{Chapter Summary}

The 2015 prize recognized the mechanistic understanding of the three repair systems that protect the genetic code: base excision repair by Lindahl, mismatch repair by Modrich, and nucleotide excision repair by Sancar. Together they explained how cells correct the damage that would otherwise lead to mutations and disease.
